\section{结论}
\label{sec:conclusion}

本文通过有界迭代阶乘的纵向案例，分析了 SysY 程序从源文本至 RV64 ELF 进程映像的表示演化。预处理确定解析器接收的记号，前端建立语法结构、名称绑定和类型；LLVM IR 将赋值、循环与调用显式化为内存效应、控制流和 SSA 数据依赖；RV64 后端在 ISA 与 psABI 约束下完成指令和寄存器映射；汇编与链接阶段通过符号、重定位和节布局确定地址关系；装载器最终依据程序头建立具有权限与入口状态的进程映像。

RV64 目标下的编译产物分析显示，\code{-O2} 相比 \code{-O0} 消除了案例 IR 中的 \code{alloca/load/store}，并将 IR 指令数由 56 降至 30、汇编近似指令行由 77 降至 45。参考 SysY、LLVM IR 和 RV64 汇编在六个边界输入上均产生一致输出和退出状态。这些结果共同支持跨层语义映射的具体解释，同时将结论限定于给定工具链、目标环境和测试范围。

各表示并非仅按抽象程度线性减少信息。AST 增加语言绑定与类型，SSA 增加定义--使用关系，对象文件增加符号与重定位，最终 ELF 增加可装载段和入口配置。MLIR 的渐进式降低把这种信息保留期限扩展为显式的体系结构选择。编译系统的跨阶段设计因而可以表述为：在保持可观察语义的条件下，为每个后续变换选择足够且可验证的表示信息。
